\section{Related Work}
\label{sec:related}

\paragraph{LLMs for causal reasoning and graph discovery.} Recent work has used language models for structured causal-discovery behavior rather than only pairwise causal judgments. \citet{lampinen2023passive} study whether agents and language models can acquire active causal strategies from passive experience; \citet{jiralerspong2024bfs} use LLMs to propose full causal graphs through a breadth-first query schedule; \citet{abdulaal2024cma} combine metadata- and data-driven reasoning to recover causal structure; and \citet{vashishtha2023causalorder} use LLM-derived causal orders as a discovery prior for downstream effect estimation. These are method neighbors to ACDB: they treat causal structure recovery as an LLM-facing problem, but they do not provide a shared benchmark that places end-to-end agents, hybrid policies, and classical active baselines under the same intervention protocol and scoring contract.

\paragraph{Benchmarks for LLM causal discovery.} Existing benchmark work splits between static causal reasoning and interactive graph discovery. \citet{jin2024corr2cause} tests whether models can infer causal direction from correlational statements alone, and \citet{chen2025autobench} evaluates iterative causal-graph identification through interventions and is the closest benchmark competitor to ACDB. \citet{havrilla2025igda} studies the complementary setting in which an agent uses semantic metadata and edge experiments rather than anonymized observational and interventional samples. \citet{realizingllmcausal2025} argues that current LLM causal-reasoning benchmarks are vulnerable to pretraining contamination and calls for science-grounded, novel evaluation protocols. ACDB sits closest to this benchmark line: it keeps the task synthetic and controlled, anonymizes variable identities to block semantic retrieval, exposes a fixed observe--intervene--submit runtime, and scores submitted graphs with a layered contract rather than a single success criterion.
